% ============================================================================
%  Adaptive Leaf Granularity for Hierarchical Retrieval
%  Working draft. Branch: ckraptor.  Status as of 2026-06-18:
%    H1 (size matters & heterogeneous) PASS at n=50; H2 (cheap density features
%    predict optimal size) FAIL at n=50 -> reported as an honest negative result.
%  Self-contained (no external .bib) so it compiles with a plain pdflatex run.
% ============================================================================
\documentclass[11pt]{article}

\usepackage[utf8]{inputenc}
\usepackage[T1]{fontenc}
\usepackage{lmodern}
\usepackage[margin=1in]{geometry}
\usepackage{amsmath,amssymb}
\usepackage{booktabs}
\usepackage{graphicx}
\usepackage{xcolor}
\usepackage{url}
\usepackage[hidelinks]{hyperref}

\newcommand{\todo}[1]{{\color{red}\textbf{[TODO: #1]}}}

\title{\textbf{When Does Adaptive Leaf Granularity Help Hierarchical
Retrieval?}\\[2pt]
\large A Training-Free Harness and a Negative Result for Cheap Density Features}

\author{Shagun Gupta\\\small Working draft --- \today}
\date{}

\begin{document}
\maketitle

\begin{abstract}
RAPTOR-style hierarchical retrieval builds a tree by recursively clustering and
summarizing \emph{leaf chunks}, but inherits a naive fixed-size leaf chunker that
ignores how documents differ in information density. We ask whether per-document
leaf granularity can be set \emph{cheaply} and \emph{without training a model},
and whether doing so beats a single well-tuned fixed leaf size inside an
otherwise-frozen pipeline. We make two contributions of method and one of
result. (i) We build a reusable, \emph{no-LLM} evaluation harness for leaf
granularity in hierarchical RAG: an oracle leaf-size sweep scored by a
gold-evidence token-coverage proxy. (ii) We frame chunking control as two cheap
deterministic axes---a \emph{structure score} that decides \emph{where} to cut
and a \emph{density score} that decides \emph{how big} each leaf is. (iii) On
QASPER (50 papers, 171 answerable questions) we find that per-document optimal
leaf size is genuinely heterogeneous, leaving real headroom---a per-document
oracle beats the best single fixed size by $+9.8\%$ relative gold-evidence
coverage (\textbf{H1}). However, cheap deterministic density features
\emph{fail} to capture that headroom: used to size leaves, the
selected-feature map beats a tuned global fixed size in only $7/50$ random
train/test splits (mean $0.704$ vs.\ $0.729$), and a learned regressor on the
same features also loses (\textbf{H2}). We report this as an honest negative
result: the common ``smaller chunks for denser text'' heuristic is
\emph{not supported} at the hierarchical-leaf level; a single tuned global leaf
size is a strong baseline. A QuALITY control confirms the structure axis does no
harm on unstructured prose (structure-routing rate $0$).
\end{abstract}

% ----------------------------------------------------------------------------
\section{Introduction}
\label{sec:intro}

Chunking is a known lever for retrieval-augmented generation (RAG), but its
interaction with \emph{hierarchical} retrieval is under-studied. RAPTOR
\cite{sarthi2024raptor} is agnostic to leaf chunking: it splits a document into
fixed-size leaves ($\approx 100$ tokens), then recursively clusters and
LLM-summarizes them into a tree that is queried by collapsed-tree retrieval. The
leaf chunker is therefore a clean, isolatable variable inside an otherwise-frozen
pipeline---and a fixed leaf size ignores that documents differ in information
density.

The \emph{broad} idea ``adapt chunk size to content'' is not novel. Our specific
question is narrower and, to our knowledge, unoccupied: can per-document leaf
granularity be set by a \emph{cheap, closed-form, training-free} score, and does
that beat a single tuned fixed size for hierarchical-RAG leaves? We answer it with
a deliberately strong methodology---an oracle upper bound, a tuned fixed-size
lower bound, and a learned regressor as a same-features ceiling---so that a
negative answer is informative rather than an artifact of a weak method.

\paragraph{Contributions.}
\begin{enumerate}
  \item A reusable \textbf{training-free, no-LLM harness} (Section~\ref{sec:harness})
  to study leaf granularity in hierarchical RAG: an oracle leaf-size sweep scored
  by a gold-evidence token-coverage proxy; resumable and free-tier-friendly.
  \item A \textbf{two-axis decomposition} of chunking control
  (Section~\ref{sec:method}): a structure score (\emph{where} to cut) and a
  density score (\emph{how big}), each deterministic and LLM-free.
  \item An empirical \textbf{characterization and refutation}
  (Sections~\ref{sec:h1}--\ref{sec:h2}): per-document optimal leaf size is
  heterogeneous with real oracle headroom (\textbf{H1}), but cheap density
  features---and a learned regressor over them---fail to capture it (\textbf{H2}).
  This is a first clean formalization and refutation of the common ``denser
  $\Rightarrow$ smaller chunks'' heuristic for hierarchical-RAG leaves.
\end{enumerate}

% ----------------------------------------------------------------------------
\section{Related Work}
\label{sec:related}

\paragraph{Hierarchical retrieval.} RAPTOR \cite{sarthi2024raptor} is our frozen
backbone; we vary only its leaf chunker.

\paragraph{Adaptive / learned chunk granularity.} Mix-of-Granularity
\cite{mog2024} \emph{learns} a soft router over granularities (a trained model)---
the closest prior art and exactly the ``use a model'' route we avoid as a
contribution. HiChunk \cite{hichunk2025} performs hierarchical structuring with a
\emph{fine-tuned} LLM and query-adaptive granularity. ``Rethinking Chunk Size for
Long-Document Retrieval'' \cite{rethink2025} analyzes chunk size but proposes no
per-document policy. Ekimetrics Adaptive Chunking \cite{ekimetrics} selects a
discrete chunking \emph{method} per document with a classifier over flat
retrieval; it does not set a continuous \emph{granularity} and trains a selector.

\paragraph{Our niche.} A \emph{cheap, closed-form, training-free} score that sets
\emph{per-section leaf granularity} in a \emph{hierarchical tree}, validated
against an oracle. ``Training-free'' separates us from the learned MoG/HiChunk;
the continuous size, the hierarchical-leaf target, and oracle calibration
separate us from Ekimetrics. The two-axis decomposition (structure$\to$boundary,
density$\to$size) is, to our knowledge, not stated cleanly in prior work.
\emph{This paper reports that the density axis, in its cheap training-free form,
does not pay off}---which is itself the result.

% ----------------------------------------------------------------------------
\section{A Training-Free Harness for Leaf Granularity}
\label{sec:harness}

Measuring ``which leaf size is best for this document'' end-to-end requires an
LLM per (document, size) cell---prohibitive on a free-tier budget. We instead
use a cheap retrieval proxy that needs \emph{no} LLM.

\paragraph{Oracle leaf-size sweep.} For each document $d$ and each candidate leaf
size $S \in \{50,100,150,200,300,400\}$ tokens, we split $d$ into $S$-token
leaves, embed them with a fixed SBERT model
(\texttt{multi-qa-mpnet-base-cos-v1}), index with FAISS, and retrieve under a
fixed $2000$-token budget for every question of $d$.

\paragraph{Gold-evidence coverage proxy.} We score retrieval by
\emph{gold-evidence coverage}: the fraction of a question's gold-evidence
paragraphs whose tokens are recalled (token-recall $\ge 0.8$) within the
retrieved context. This replaces exact-paragraph evidence-$F_1$, which is
ill-defined when chunk boundaries differ from gold-paragraph boundaries. We
define the per-document optimal size as
$\hat S(d) = \arg\max_S \mathrm{cov}(d,S)$.

The sweep is deterministic, resumable, and Drive-cached; the QASPER run produces
$50 \times 6 = 300$ \texttt{(doc,size)} records with no LLM call. This is what
makes a rigorous negative result affordable.

% ----------------------------------------------------------------------------
\section{Method: Two-Axis Chunking Control}
\label{sec:method}

We decompose leaf chunking into two cheap deterministic axes; everything below is
LLM-free at inference.

\paragraph{Axis 1 --- structure score (where to cut).} A deterministic
$\texttt{structure\_score}(d)\in[0,1]$ from surface features (numbered / markdown
/ ALL-CAPS / keyword headings, a generic Title-Case heading detector, TOC-like
regions, inter-heading spacing regularity). A document scoring above a threshold
$\tau$ is routed to a structure-aware chunker (boundaries at headings); otherwise
it falls back to token chunking. This is the Adaptive Hybrid Chunking (AHC) route.

\paragraph{Axis 2 --- density score (how big).} A deterministic
$\texttt{density\_score}(d) = (\rho \in [0,1], \mathbf{f})$ from six surface
features $\mathbf{f}$: mean sentence length, type--token ratio, numeral/symbol
density, mean word length, list-marker density, and non-stopword ratio. A
calibrated map $S = a + b\,f$ turns one chosen feature into a leaf-token budget.

\paragraph{Calibration, not learning.} Rather than assume a direction (``denser
$\Rightarrow$ smaller''), we \emph{calibrate} the density$\to$size map against the
oracle sweep of Section~\ref{sec:harness}. On a training split we rank features by
$|\text{Spearman}|$ against $\hat S(d)$, select the single strongest feature, and
fit a two-parameter closed form $S = a + b\,f$ (no test-set tuning). We compare
this training-free map against (a) the best global fixed size, (b) a learned
$6$-feature linear regressor (same features, a ceiling for ``what these features
can do''), and (c) the per-document oracle (an upper bound).

% ----------------------------------------------------------------------------
\section{Experimental Setup}
\label{sec:setup}

\textbf{Data.} QASPER (scientific papers); $50$ papers / $171$ answerable
(gold-evidence) questions. \textbf{Embedder/reader.}
\texttt{multi-qa-mpnet-base-cos-v1} SBERT; \texttt{gpt-oss-120b} (free tier) where
an LLM is used. \textbf{Proxy metric.} gold-evidence token-coverage
($\ge 0.8$ recall). \textbf{Protocol.} $50$ random $70/30$ train/test splits;
we report mean$\pm$std TEST coverage per policy, the win-rate of the selected
feature over the best global fixed size, and feature-selection stability.
\textbf{Controls.} Only the chunker varies; embeddings, clustering, retrieval
budget, and reader are frozen across arms. Results are within-pipeline
(arm-vs-arm), not leaderboard numbers.

% ----------------------------------------------------------------------------
\section{H1: Leaf Size Matters and the Optimum is Heterogeneous (PASS)}
\label{sec:h1}

The best single fixed size is $200$ tokens at coverage $0.7362$; the per-document
oracle reaches $0.8080$---an absolute gap of $0.0718$, or $\mathbf{+9.8\%}$
\textbf{relative} gold-evidence coverage left on the table by any fixed size. The
per-document optimum is spread across the whole grid
(Table~\ref{tab:hetero}), with mass at every size. There \emph{is} adaptive
headroom; H1 is the premise, not the contribution.

\begin{table}[t]
\centering
\caption{H1: distribution of per-document optimal leaf size over 50 QASPER papers.
The optimum spans the full grid---adaptive headroom exists.}
\label{tab:hetero}
\begin{tabular}{lcccccc}
\toprule
optimal size (tok) & 50 & 100 & 150 & 200 & 300 & 400 \\
\midrule
\# documents & 17 & 11 & 6 & 8 & 3 & 5 \\
\bottomrule
\end{tabular}
\end{table}

% ----------------------------------------------------------------------------
\section{H2: Cheap Density Features Do Not Capture the Headroom (FAIL)}
\label{sec:h2}

\paragraph{Feature signal is weak.} Table~\ref{tab:spearman} shows the per-feature
Spearman correlation against $\hat S(d)$ at $n{=}50$. The strongest single feature
is \texttt{type\_token\_ratio} at only $-0.379$; every other feature has
$|\rho|\le 0.16$, and the equal-weight composite $\rho$ is $-0.147$ (the features
are oppositely signed, so a uniform mean cancels them). The features that looked
strong on an 8-document preview ($\texttt{mean\_word\_len}\;{+}0.79$,
$\texttt{nonstopword\_ratio}\;{+}0.66$) collapsed at scale---small-$N$
overfitting.

\begin{table}[t]
\centering
\caption{H2 feature ablation: Spearman of each density signal vs.\ per-document
optimal size ($n{=}50$), with the $n{=}8$ preview for contrast. The preview's
strong features collapse at scale.}
\label{tab:spearman}
\begin{tabular}{lrr}
\toprule
signal & $n{=}50$ & ($n{=}8$ preview) \\
\midrule
\texttt{type\_token\_ratio}          & $-0.379$ & $-0.602$ \\
\texttt{nonstopword\_ratio}          & $+0.158$ & $+0.663$ \\
\texttt{numeral\_symbol\_density}    & $+0.153$ & $-0.258$ \\
\texttt{mean\_word\_len\_chars\_norm}& $+0.108$ & $+0.786$ \\
\texttt{mean\_sentence\_len\_norm}   & $-0.080$ & $+0.025$ \\
\texttt{list\_marker\_density}       & $+0.061$ & $\;\;0.000$ \\
$\rho$ (equal-weight composite)      & $-0.147$ & $+0.037$ \\
\bottomrule
\end{tabular}
\end{table}

\paragraph{Every adaptive policy loses to a tuned fixed size.}
Table~\ref{tab:robust} reports the headline statistic: held-out coverage over
$50$ random $70/30$ splits. The selected-feature map beats the best global fixed
size in only $7/50$ splits ($14\%$) and is below it on average
($0.704$ vs.\ $0.729$). The selected feature \emph{is} stable
(\texttt{type\_token\_ratio} chosen in $47/50$ splits)---so this is a
\emph{stably weak} predictor, not noise. Crucially, the \emph{learned} regressor
also loses ($0.707$): the bottleneck is the feature set, not the closed-form map.
None of the predictors approaches the oracle ($0.806$); essentially all of the
H1 headroom is left unrealized.

\begin{table}[t]
\centering
\caption{H2 held-out robustness over 50 random $70/30$ splits (QASPER, $n{=}50$).
Every adaptive/learned policy lands below the tuned global fixed size, and far
below the oracle.}
\label{tab:robust}
\begin{tabular}{lrr}
\toprule
policy (TEST docs) & mean & std \\
\midrule
\textbf{best global fixed (200 tok)}       & \textbf{0.7291} & 0.0433 \\
training-free: composite $\rho \to$ size    & 0.7212 & 0.0447 \\
training-free: selected feature $\to$ size  & 0.7037 & 0.0494 \\
learned regressor (6 features)              & 0.7071 & 0.0415 \\
per-document oracle (ceiling)               & 0.8061 & 0.0347 \\
\bottomrule
\end{tabular}
\end{table}

\paragraph{Verdict.} Against a pre-registered gate---(i) high win-rate over fixed,
(ii) stable selected feature, (iii) approaches learned \emph{and} oracle---only
(ii) is met. \textbf{H2 fails.} A stable predictor that reliably loses to a fixed
baseline is just stably weak, and because the learned ceiling fails too, no model
sophistication on these features converts the headroom into a win.

% ----------------------------------------------------------------------------
\section{Structure-Axis Control: No Harm on Unstructured Text}
\label{sec:control}

On QuALITY (plain narrative prose; multiple-choice), AHC and the token baseline
are a statistical tie ($0.400$ vs.\ $0.381$ accuracy over $105$ paired questions;
overlapping $95\%$ CIs). The cause is mechanical and intended: the
structure-routing rate is $0.0000$---no document has exploitable structure, so AHC
falls back to token chunking everywhere and runs the same chunker as the baseline.
QuALITY thus serves as an \emph{unstructured control}: the structure axis does no
harm when there is nothing to exploit. The structure path is genuinely exercised
only on structured corpora (QASPER); that end-to-end run is in progress.

% ----------------------------------------------------------------------------
\section{Discussion}
\label{sec:discussion}

\paragraph{What we refute.} On QASPER, per-document leaf-size heterogeneity
creates real oracle headroom ($+9.8\%$ rel.), but cheap, training-free density
features---and a learned regressor on them---do not predict per-document optimal
leaf size well enough to beat a single tuned global fixed size. The popular
``smaller chunks for denser text'' heuristic is not supported at this granularity
for hierarchical-RAG leaves.

\paragraph{Practical takeaway.} Tune \emph{one} global leaf size ($\approx 200$
tokens here); it is a strong, hard-to-beat baseline. Per-document density-based
leaf sizing is not worth its complexity without a better signal than closed-form
surface density.

\paragraph{Where a win could still live.} The headroom is real and unclaimed: the
oracle leaves $\approx 0.10$ coverage above every predictor we tried. A richer,
\emph{query-aware} or learned multi-signal predictor (in the spirit of
\cite{mog2024,hichunk2025}) could still capture it. We refute only the cheap,
closed-form, training-free density route---the project's intended contribution.

% ----------------------------------------------------------------------------
\section{Threats to Validity}
\label{sec:threats}

\textbf{Proxy gap.} Gold-evidence token-coverage (retrieval) is not end-task QA;
results are arm-vs-arm. \textbf{Scale / corpus.} $N{=}50$, single corpus (QASPER);
CIs are wide (std $\approx 0.04$). The negative conclusion is robust \emph{for the
cheap-density route} precisely because the learned ceiling on the same features
also fails, but it does not generalize to other document types or to richer
signals. \textbf{Free-tier nondeterminism.} Clustering and LLM steps are not
bit-exact; we report variance and cache all LLM calls. \textbf{Fixed
embedder/reader.} Within-pipeline numbers only.

% ----------------------------------------------------------------------------
\section{Conclusion}
\label{sec:conclusion}

We built a training-free, no-LLM harness for studying leaf granularity in
hierarchical RAG and used it to show that, while per-document optimal leaf size is
heterogeneous and leaves real oracle headroom (H1), cheap deterministic density
features cannot capture it (H2)---a tuned global fixed size wins, and a learned
regressor on the same features does not change that. We report this as an honest,
reproducible refutation of a common-but-unvalidated chunking heuristic, and we
release the harness so that richer per-document granularity predictors can be
evaluated against the same oracle.

% ----------------------------------------------------------------------------
\begin{thebibliography}{9}
\bibitem{sarthi2024raptor}
P.~Sarthi, S.~Abdullah, A.~Tuli, S.~Khanna, A.~Goldie, and C.~D.~Manning.
\newblock RAPTOR: Recursive Abstractive Processing for Tree-Organized Retrieval.
\newblock \emph{ICLR}, 2024. arXiv:2401.18059.

\bibitem{mog2024}
Z.~Zhong et al.
\newblock Mix-of-Granularity: Optimize the Chunking Granularity for
Retrieval-Augmented Generation.
\newblock arXiv:2406.00456, 2024.

\bibitem{hichunk2025}
HiChunk: Hierarchical Document Chunking with Auto-Merge.
\newblock arXiv:2509.11552, 2025.

\bibitem{rethink2025}
Rethinking Chunk Size for Long-Document Retrieval.
\newblock arXiv:2505.21700, 2025.

\bibitem{ekimetrics}
Ekimetrics.
\newblock Adaptive Chunking (chunking-method selection).
\newblock \url{https://github.com/ekimetrics/adaptive-chunking}.
\end{thebibliography}

\end{document}
